\documentclass{article}
\usepackage{graphicx}
\usepackage{minted}
\usepackage{titling}
\usepackage[left=3cm,right=3cm,top=2cm,bottom=2cm]{geometry}
\usepackage{booktabs}
\usepackage{multirow}
\usepackage{enumitem}

\setlength{\droptitle}{-6 em}
\pagenumbering{gobble}
\renewcommand{\arraystretch}{1.3}

\title{Circuitos Digitais II: Aulas 2-3}
\author{Luis Felipe Ferreira Soares}
\date{13/11/2025}

\begin{document}

\maketitle

\section{Flip-Flops}
\subsection{}
\begin{table}[h!]
\centering
\label{tab:not_clk}
\begin{tabular}{|c|c|c|c|c|c|c|}
    \hline
    Circuito & CLK & D & R & S & Q & $\overline{Q}$\\
    \hline
    \multirow{6}{*}{\includegraphics[width=0.2\textwidth]{1_1.png}}
    & $\downarrow$ & 0 & 0 & 0 & 0 & 1\\
    \cline{2-7}
    & $\downarrow$ & 1 & 0 & 0 & 1 & 0\\
    \cline{2-7}
    & $\uparrow$ & x & 0 & 0 & Q & $\overline{Q}$\\
    \cline{2-7}
    & x & x & 1 & 0 & 0 & 1\\
    \cline{2-7}
    & x & x & 0 & 1 & 1 & 0\\
    \cline{2-7}
    & x & x & 1 & 1 & 1 & 1\\
    \hline
\end{tabular}
\end{table}

\subsection{}
\begin{table}[h!]
\centering
\label{tab:not_reset}
\begin{tabular}{|c|c|c|c|c|c|c|}
    \hline
    Circuito & CLK & D & R & S & Q & $\overline{Q}$\\
    \hline
    \multirow{6}{*}{\includegraphics[width=0.2\textwidth]{1_2.png}}
    & $\uparrow$ & 0 & 1 & 0 & 0 & 1\\
    \cline{2-7}
    & $\uparrow$ & 1 & 1 & 0 & 1 & 0\\
    \cline{2-7}
    & $\downarrow$ & x & 1 & 0 & Q & $\overline{Q}$\\
    \cline{2-7}
    & x & x & 0 & 0 & 0 & 1\\
    \cline{2-7}
    & x & x & 1 & 1 & 1 & 0\\
    \cline{2-7}
    & x & x & 0 & 1 & 1 & 1\\
    \hline
\end{tabular}
\end{table}

\subsection{}
\begin{table}[h!]
\centering
\label{tab:not_reset_clk}
\begin{tabular}{|c|c|c|c|c|c|c|}
    \hline
    Circuito & CLK & D & R & S & Q & $\overline{Q}$\\
    \hline
    \multirow{6}{*}{\includegraphics[width=0.2\textwidth]{1_2.png}}
    & $\downarrow$ & 0 & 1 & 0 & 0 & 1\\
    \cline{2-7}
    & $\downarrow$ & 1 & 1 & 0 & 1 & 0\\
    \cline{2-7}
    & $\uparrow$ & x & 1 & 0 & Q & $\overline{Q}$\\
    \cline{2-7}
    & x & x & 0 & 0 & 0 & 1\\
    \cline{2-7}
    & x & x & 1 & 1 & 1 & 0\\
    \cline{2-7}
    & x & x & 0 & 1 & 1 & 1\\
    \hline
\end{tabular}
\end{table}

\section{Quais as vantagens dos Latches? Quais são as suas desvantagens?}
\begin{itemize}[leftmargin=*]
    \item \textbf{Vantagens}: Costumam ser mais rápidos que flip-flops de borda, pois podem mudar de estado a qualquer momento que o sinal de enable esteja em alto, sem necessidade de aguardar pela subida de borda como nos flip-flops, característica que também é útil em aplicações de controle de memória. Em termos físicos, demandam menos portas lógicas que os flip-flops de borda, portanto ocupam menos espaço e consomem menos energia.
    \item \textbf{Desvantagens}: A transparência em relação a entrada enquanto o enable está ativo também pode ser uma desvantagem, pois em sistemas de alta sensibilidade isto pode ocasionar glitches e condições de corrida inesperadas, mesmo quando o enable está diretamente conectado a um clock.
\end{itemize}

\section{Realize uma pesquisa e descreva em quais situações (ou aplicações) podem/devem ser usados Latches.}
São utilizados em Static Random Access Memory (SRAMs) devido a sua velocidade, buffers e registradores intermediários/temporários pelo mesmo motivo, amostragem/inserção de dados de modo assíncrono e em sistemas baseados em \textbf{Latching Logic}, que utiliza as principais vantagens dos latches para otimizar o tempo entre os estágios de uma pipeline, permitindo que os dados propaguem por mais tempo, otimizando o caminho crítico.

\section{Construa em Verilog ou VHDL um multiplexador de 4 entradas e uma saída com largura de bits parametrizável. Use sentenças IF – ELSIF para produzir a presença de Latches. Realize a síntese lógica e analise o esquemático RTL. A ferramenta inferiu Latches? Quantos Latches foram inferidos? Realize um testbench e verifique o funcionamento através de simulação.}
\subsection{Código em Verilog}
\begin{minted}{verilog}
module mux4to1_p #(parameter N = 4) (in3, in2, in1, in0, sel, out);
input [N-1:0] in3, in2, in1, in0;
input [1:0] sel;
output reg [N-1:0] out;

always @* begin
    if (sel == 2'b00) out = in0;
    else if (sel == 2'b01) out = in1;
    else if (sel == 2'b10) out = in2;
    else if (sel == 2'b11) out = in3;
end

endmodule
\end{minted}

\subsection{Foram inferidos 4 latches}
\includegraphics[width=1\textwidth]{2_3.png}

\subsection{Testbench}
\begin{minted}{verilog}
module mux4to1_p_tb;
parameter N = 4;
reg [N-1:0] D3, D2, D1, D0;
reg [1:0] sel;
wire [N-1:0] out;

mux4to1_p#(N) dut (D3, D2, D1, D0, sel, out);

always #1 sel[0] = !sel[0];
always #2 sel[1] = !sel[1];

initial begin
    $display("  D3  |  D2  |  D1  |  D0  | sel | out ");
    $monitor(" %4b | %4b | %4b | %4b | %3b | %b", D3, D2, D1, D0, sel, out);
    sel = 0;
    D3 = 4'b1111; D2 = 4'b0111; D1 = 4'b0011; D0 = 4'b0001;
    #4 $finish(0);
end

endmodule
\end{minted}

\subsection{Simulação}
\includegraphics[width=1\textwidth]{2_3-wave.png}

\newpage

\section{Repita o exercício anterior usando diretivas CASE.}
\subsection{Código em Verilog}
\begin{minted}{verilog}
module mux_v2 #(parameter N = 4) (in3, in2, in1, in0, sel, out);
input [N-1:0] in3, in2, in1, in0;
input [1:0] sel;
output reg [N-1:0] out;

always @* begin
    case (sel)
        2'b00: out = in0;
        2'b01: out = in1;
        2'b10: out = in2;
        2'b11: out = in3;
    endcase
end

endmodule
\end{minted}

\subsection{Não foram inferidos latches}
\includegraphics[width=1\textwidth]{2_4.png}

\newpage

\subsection{Testbench}
\begin{minted}{verilog}
module mux_v2_tb;
parameter N = 4;
reg [N-1:0] D3, D2, D1, D0;
reg [1:0] sel;
wire [N-1:0] out;

mux_v2#(N) dut (D3, D2, D1, D0, sel, out);

always #1 sel[0] = !sel[0];
always #2 sel[1] = !sel[1];

initial begin
    $display("  D3  |  D2  |  D1  |  D0  | sel | out ");
    $monitor(" %4b | %4b | %4b | %4b | %3b | %b", D3, D2, D1, D0, sel, out);
    sel = 0;
    D3 = 4'b1111; D2 = 4'b0111; D1 = 4'b0011; D0 = 4'b0001;
    #4 $finish(0);
end

endmodule
\end{minted}

\subsection{Simulação}
\includegraphics[width=1\textwidth]{2_4-wave.png}

\end{document}